\documentclass[11pt]{article}
\usepackage[a4paper,margin=21mm]{geometry}
\usepackage[no-math]{fontspec}
\setmainfont{Latin Modern Roman}
\newfontfamily\CJKfont{Noto Serif CJK SC}
\XeTeXlinebreaklocale "zh"
\XeTeXlinebreakskip=0pt plus 1pt
\usepackage{amsmath,amssymb,amsthm,amscd,graphicx,color}
\usepackage{xurl}
\usepackage[unicode,hidelinks]{hyperref}
\allowdisplaybreaks
\setlength\emergencystretch{3em}
\numberwithin{equation}{section}

\usepackage{mathrsfs}
\usepackage{enumitem}

\numberwithin{equation}{section}

\newtheorem{Theorem}{Theorem}[section]
\newtheorem*{Theorem*}{Theorem}
\newtheorem{Corollary}[Theorem]{Corollary}
\newtheorem{Lemma}[Theorem]{Lemma}
\newtheorem{Proposition}[Theorem]{Proposition}
\newtheorem{obs}[Theorem]{Observation}
 { \theoremstyle{definition}
\newtheorem{Definition}[Theorem]{Definition}
\newtheorem{Note}[Theorem]{Note}
\newtheorem{Example}[Theorem]{Example}
\newtheorem{Remark}[Theorem]{Remark}
\newtheorem{exas}[Theorem]{Examples}
\newtheorem{exa}[Theorem]{Example}}

\DeclareMathOperator{\im}{im}
\DeclareMathOperator{\Ad}{Ad}
\DeclareMathOperator{\ad}{ad}
\DeclareMathOperator{\Area}{Area}
\DeclareMathOperator{\Cut}{Cut}
\DeclareMathOperator{\Def}{Def}
\DeclareMathOperator{\diver}{div}
\DeclareMathOperator{\dom}{Dom}
\DeclareMathOperator{\ess}{ess}
\DeclareMathOperator{\essran}{essran}
\DeclareMathOperator{\Hess}{Hess}
\DeclareMathOperator{\id}{id}
\DeclareMathOperator{\Lie}{Lie}
\DeclareMathOperator{\loc}{loc}
\DeclareMathOperator{\rad}{rad}
\DeclareMathOperator{\ran}{Ran}
\DeclareMathOperator{\Ray}{Ray}
\DeclareMathOperator{\Ric}{Ric}
\DeclareMathOperator{\supp}{supp}
\DeclareMathOperator{\tr}{tr}

\def\p{\mathrm{pp}}
\def\c{\mathrm c}
\def\d{\mathrm d}
\def\ac{\mathrm{ac}}
\def\sc{\mathrm{sc}}
\def\s{\mathrm s}
\def\C{\mathbb C}
\def\F{\mathbb F}
\def\N{\mathbb N}
\def\R{\mathbb R}
\def\H{\mathbb H}
\def\A{\mathcal A}
\def\T{\mathcal T}
\def\ve{\varepsilon}
\def\vf{\varphi}

\makeatletter\newlabel{secspec}{{A}{}{Original source locator}{}{}}\makeatother
\begin{document}
\begin{center}\Large Absolutely continuous Laplace spectrum on negatively curved Hadamard manifolds asymptotically harmonic at one ideal point\end{center}
\noindent\textbf{Stable ID:} 1147\quad\textbf{Author:} Werner Ballmann; Mayukh Mukherjee; Panagiotis Polymerakis\par
\noindent\textbf{Work:} On the Spectrum of Certain Hadamard Manifolds\par
\noindent\textbf{Source:} \url{https://sigma-journal.com/2023/050/}\par
\noindent\textbf{Version:} Official SIGMA 19 (2023) article 050, DOI 10.3842/SIGMA.2023.050; journal PDF and native TeX acquired 2026-09-30, exact bytes pinned by SHA256\par
\noindent\textbf{Rights:} CC BY 4.0. Authors retain copyright. \url{https://creativecommons.org/licenses/by/4.0/}. Reuse and typesetting adaptation; original language and mathematical content preserved.\par
\noindent\textbf{Difficulty (prior selection preserved):} {\CJKfont H2: The proof uses a Busemann vector field to derive a positive commutator estimate controlling derivatives tangent to horospheres. It converts this into local smoothness for the resolvent and excludes every singular spectral component through distributional constancy along infinite-volume horospheres. For the pinched-curvature case it builds cutoff pullbacks of one-dimensional approximate eigenfunctions, proving the entire half-line is present in the spectrum.}\par
\medskip\noindent\textbf{Editorial boundary note (not author text):} Complete original Theorem1.2 including both absolute continuity and the stronger full-half-line spectrum under the actual extra pinching/curvature-derivative hypotheses. Hadamard/asymptotically harmonic/Laplacian-domain and spectral-measure conventions, Busemann and Cheeger/Hessian observations with full current proofs, full Rellich identity proof and integrated positive-commutator chain, actual smooth-resolvent criteria, and complete local tangent-derivative/core/distributional argument are retained. The entire current one-dimensional conjugation, flow-invariant cutoff/Jacobi estimate, published geometric Laplacian-bound application and every norm/error/Weyl sequence calculation are included; the author expressly subsumes M by the trivial-group quotient at source931. Quotient framing and its current supporting argument are preserved without extra problem counts. Named resolvent smoothness criteria and the explicitly stated prior geometric cutoff estimate are precise established prerequisites, not omitted new current constructions. One expressly announced second main spectral outcome; homogeneous, boundary-domain and product claims are not separately selected. Literal author source664 L2(D) wording in the otherwise M-based proof and all norm/error signs remain unchanged, without editorial identification or verification.\par
\noindent\textbf{External original reference secspec:} Original spectral-theory appendix, source1054–1332. Mentioned in the retained introductory spectral conventions for the separate product discussion. The selected proof states its actual smooth-resolvent criteria and uses the precise Reed–Simon/Xavier citations; it does not defer its new Busemann/cutoff construction to this appendix.\par
\medskip\hrule\medskip\noindent\textbf{Author text begins below.}\par
\setcounter{section}{1}
% Original source lines 124--164
For a complete and connected Riemannian manifold $M$, we view its Laplacian $\Delta$
as an unbounded and symmetric operator with domain $C^\infty_{\rm c}(M) \subseteq L^2(M)$.
The closure of $\Delta$, which we also denote by $\Delta$, is a self-adjoint operator in $L^2(M)$.
Its spectrum as a subset of $\R$ will be denoted by $\sigma(M)$.

Functional analysis of self-adjoint operators yields the orthogonal decomposition
\begin{align*}%\label{deco}
L^2(M) = H_{\p}(M) \oplus H_{\ac}(M) \oplus H_{\sc}(M)
\end{align*}
of $L^2(M)$ into $\Delta$-invariant subspaces such that
\begin{enumerate}\itemsep=0pt
\item[(1)] $H_{\p}(M)$ is spanned by eigenfunctions;
\item[(2)] $H_{\ac}(M)$ consists of functions with absolutely continuous spectral measure;
\item[(3)] $H_{\sc}(M)$ consists of functions with singularly continuous spectral measure.
\end{enumerate}
We present some more details about these notions and decompositions in Appendix~\ref{secspec},
since they are needed in our discussion of spectra of Riemannian products.

There is the well known decomposition $\sigma(M)=\sigma_{\d}(M)\cup\sigma_{\ess}(M)$ of the spectrum
as a~disjoint union of \emph{discrete} and \emph{essential spectrum},
where $\lambda\in\R$ belongs to $\sigma_{\ess}(M)$ if $\Delta-\lambda$ is not a~Fredholm operator.
Then $\lambda\in\sigma_{\d}(M)$ if and only if $\lambda$ is an eigenvalue of $\Delta$
of finite multiplicity and is an isolated point of $\sigma(M)$.
Therefore, the preimage in $L^2(M)$ of $\sigma_{\d}(M)$ under the spectral projection associated to $\Delta$
is contained in $H_{\p}(M)$.
In particular, if the essential spectrum of~$\Delta$ is empty, then $H_{\ac}(M)=H_{\sc}(M)=\{0\}$.
The essential spectrum of $\Delta$ is determined by the geometry of $M$ at infinity;
see, for example,~\cite[Proposition 3.6]{BMM17} for a general formulation of this fact.
In particular, the essential spectrum vanishes if $M$ is compact.
But there are also a~number of results about the vanishing of the essential spectrum
for non-compact manifolds;
see, for example,~\cite[Theorems 1.1 and 1.2]{BLM10} and~\cite[Examples 3.7]{BMM17}.

Now the vanishing of the essential spectrum implies the vanishing of the absolutely continuous spectrum,
which is not what we are aiming for.
In fact, we will obtain conditions which imply the absolute continuity of the spectrum, $H_{\ac}(M)=L^2(M)$,
or, in other words, the vanishing of the point and the singular continuous spectrum.
As a byproduct, we will also discuss the vanishing of the point spectrum, $H_{\p}(M)=\{0\}$.
The underlying manifolds will be \emph{Hadamard manifolds}, that is,
complete and simply connected Riemannian manifolds of non-positive sectional curvature.


\setcounter{Theorem}{1}
% Original source lines 197--243
We say that a Hadamard manifold $M$ is \emph{asymptotically harmonic about a point $\xi\in M_\infty$}
if there is a number $h\ge0$ such that each horosphere in $M$ with center $\xi$ has constant mean curvature $h$.
Hyperbolic spaces $H^k_{\F}$, endowed with Fubini--Study metrics, are asymptotically harmonic about any point in $M_\infty$, where $\lambda_0=h^2/4>0$.
Asymptotically harmonic manifolds in the usual sense are asymptotically harmonic about any point in $M_\infty$.
Such manifolds arise, for example, in investigations on the rigidity of geodesic flows.
Furthermore, since horospheres are limits of spheres,
harmonic Hadamard manifolds are asymptotically harmonic about any point in $M_\infty$.
However, homogeneous Hadamard manifolds without Euclidean factor,
which are not hyperbolic spaces, are asymptotically harmonic about some, but not any point $M_\infty$.
A second main result of this article is

\begin{Theorem}\label{onlyacm}
If $K_M<0$ and $M$ is asymptotically harmonic about a point $\xi\in M_\infty$,
then the spectrum of $\Delta$ on $M$ is absolutely continuous with $\sigma(M)\subseteq\bigl[h^2/4,\infty\bigr)$,
where $h>0$ is the mean curvature of the horospheres with center $\xi$.
Moreover, if the sectional curvature of $M$ is negatively pinched and the covariant derivative of the curvature tensor of $M$ is uniformly bounded, then $\sigma(M)=\bigl[h^2/4,\infty\bigr)$.
\end{Theorem}

Say that a group of isometries of a Hadamard manifold $M$ is \emph{elementary}
if it is discrete and without torsion, fixes a point $\xi\in M_\infty$,
and such that it is of one of the following two types:
\begin{enumerate}[label=(\alph*)]\itemsep=0pt
\item\label{ax} $\Gamma$ leaves a geodesic $\gamma$ in $M$ emanating from $\xi$ invariant.
\item\label{pa} $\Gamma$ leaves Busemann functions associated to $\xi$ invariant.
\end{enumerate}
In the first case, $C=\Gamma\backslash\gamma$ is a circle and the projection $M\to\gamma$
along horospheres with center~$\xi$ descends to a Riemannian submersion $\pi\colon N\to C$,
where $N=\Gamma\backslash M$.
In the second case,
any Busemann function $b$ associated to $\xi$ is invariant under $\Gamma$
and hence pushes down to $N$ and defines a Riemannian submersion $\pi=b\colon N\to\R$.

\begin{Theorem}\label{onlyacn}
For $M$ as in Theorem~{\rm \ref{onlyacm}}, if $\Gamma$ is an elementary group of isometries of $M$ fixing~$\xi$,
then the spectrum of the quotient $N=\Gamma\backslash M$ is absolutely continuous
with $\sigma(N)\subseteq\bigl[h^2/4,\infty\bigr)$, where $h>0$ denotes the mean curvature of the fibers of $\pi$.
Moreover,
\begin{enumerate}[label={\rm (\arabic*)}]\itemsep=0pt
\item
if $N$ is of type {\rm \ref{pa}} and the fibers of $\pi$ are of finite volume or
\item
if $N$ is of type {\rm \ref{pa}}, the fibers of $\pi$ are of infinite volume,
the sectional curvature of $M$ is negatively pinched,
and the covariant derivative of the curvature tensor of $M$ is uniformly bounded,
\end{enumerate}
then $\sigma(N)=\bigl[h^2/4,\infty\bigr)$.
\end{Theorem}


% Original source lines 267--312
\section{Preliminaries}
%%%%%%%%%%%%%%%%%%%%%%%%%%%%%%%%%%%%%%%%%%%%%%

For a Hadamard manifold $M$ and points $\xi\in M_\infty$ and $x\in M$,
the Busemann function $b$ associated to $\xi$ with $b(x)=0$ is given by
\begin{align*}
	b(y) = \lim_{n\to\infty} (d(x_n,y)-d(x_n,x)),
\end{align*}
where $(x_n)$ is any sequence in $M$ converging to $\xi$.
Recall that the limit exists and that $b$ is a $C^2$ distance function on $M$
with horoballs and horospheres centered at $\xi$ as sublevel and level sets;
cf.~\cite{Ba95} for this and other facts about Hadamard manifolds.

\begin{obs}\label{chee}
If a Hadamard manifold $M$ is asymptotically harmonic about a point $\xi\in M_\infty$
or if $N=\Gamma\backslash M$ is of type {\rm \ref{pa}},
then the Cheeger constant of $M$ respectively $N$ is at least $h$,
where~$h$ is the mean curvature of the horospheres with center $\xi$.
In particular, the spectrum of~$M$ respectively $N$ is contained in $\bigl[h^2/4,\infty\bigr)$.
\end{obs}

Throughout the article,
we use the absolute value sign to denote the appropriate volumes of sets under discussion.

\begin{proof}[Proof of Observation~\ref{chee}]
Let $D$ be a compact domain in $M$ respectively $N$ with smooth boundary.
Let $b$ be a Busemann function associated to $\xi$ and $X=\nabla b$.
Then $|X|=1$ and $\diver X=h$ and hence
\begin{align*}
	h|D| = \int_D\diver X = \int_{\partial D}\langle X,\nu\rangle \le |\partial D|,
\end{align*}
where $\nu$ is the outer normal field of $D$ along $\partial D$.
The last assertion is just the Cheeger in\-equality.
\end{proof}

\begin{obs}\label{obs}
If a Hadamard manifold $M$ is asymptotically harmonic about a point $\xi\in M_\infty$
and $b$ is a Busemann function of $M$ centered at $\xi$, then
\begin{align*}
	0 \le \nabla^2b(v,v) \le h |v|^2.
\end{align*}
\end{obs}

\begin{proof}
The claim follows immediately from $h=\tr\nabla^2b$ and $\nabla^2b\ge0$.
\end{proof}

\setcounter{Theorem}{3}
% Original source lines 339--401
For a boundary condition $B$ of a domain $D$ in a Riemannian manifold $M$,
denote by $C^\infty_{c,B}(\bar{D})$ the space of smooth functions on $\bar D$ with compact support
which satisfy $B$ along the boundary~$\partial D$ of $D$.

\begin{Definition}%\label{bcg}
We say that a boundary condition $B$ for a smooth domain $D$ is \emph{non-positive}
if~$u\in C^\infty_{c,B}(\bar D)$ implies $2u\nabla u=\nabla_\nu u^2\le0$ along $\partial D$.
\end{Definition}

\begin{exas}
Dirichlet and Neumann boundary conditions are non-positive.
Robin boundary conditions $\alpha u+\beta\nabla_\nu u=0$ are non-positive if $\alpha\beta>0$.
These kinds of boundary conditions are also self-adjoint and elliptic in the usual sense.
\end{exas}

%%%%%%%%%%%%%%%%%%%%%%%%%%%%%%%%%%%%%%%%%%%%%%
\section{Vanishing of point spectrum}\label{secrell}
%%%%%%%%%%%%%%%%%%%%%%%%%%%%%%%%%%%%%%%%%%%%%%

For a $C^1$ vector field $X$ and a $C^2$ function $u$ on a Riemannian manifold $M$, we have the following identity
\begin{align}
 2\langle\nabla_{\nabla u}X,\nabla u\rangle
 = 2\langle X,\nabla u\rangle\Delta u + |\nabla u|^2\diver X + \diver\big\{2X(u)\nabla u-|\nabla u|^2X\big\}. \label{rell}
\end{align}
In the context of spectral theory, this identity was used by Donnelly--Garofalo,
see~\cite[Lemma~2.1]{DoGa92}, \cite[Proposition 3.1]{DoGa97}, and Xavier, see~\cite[identity~(1)]{Xavier1988},
but it also occurs in the earlier work of Kazdan--Warner~\cite[identity~(8.1)]{KazdanWarner74c}.
Donnelly and Garofalo attribute \eqref{rell} to Rellich~\cite{Rellich43} in the case of Euclidean spaces;
cf.~(3) respectively II on page 62 of~\cite{Rellich43}.
For the convenience of the reader, we give a~short proof of the identity.

\begin{proof}[Proof of \eqref{rell}]
Since $X(u)=\langle\nabla u, X\rangle$ and $|\nabla u|^2=\langle\nabla u,\nabla u\rangle=\nabla u(u)$, we have
\begin{align*}
 2\diver(X(u)\nabla u)
 &= 2\nabla u(X(u)) + 2X(u)\diver\nabla u \\
 &= 2X(\nabla u(u)) + 2[\nabla u,X](u) - 2X(u)\Delta u \\
 &= 2X\bigl(|\nabla u|^2\bigr) + 2(\nabla_{\nabla u}X)(u) - 2(\nabla_X\nabla u)(u) - 2X(u)\Delta u \\
 &= 2X\bigl(|\nabla u|^2\bigr) + 2\langle\nabla_{\nabla u}X,\nabla u\rangle - 2\langle\nabla_X\nabla u,\nabla u\rangle - 2X(u)\Delta u \\
 &= X\bigl(|\nabla u|^2\bigr) + 2\langle\nabla_{\nabla u}X,\nabla u\rangle - 2X(u)\Delta u \\
 &= \diver\bigl(|\nabla u|^2X\bigr) - |\nabla u|^2\diver X + 2\langle\nabla_{\nabla u}X,\nabla u\rangle - 2X(u)\Delta u,
\end{align*}
which is the assertion.
\end{proof}

\begin{Corollary}%\label{codoga}
If $\Delta u=\varphi u$ for a $C^1$ function $\varphi$ on $M$, then
\begin{gather}
 2\langle\nabla_{\nabla u}X,\nabla u\rangle
= \varphi X(u^2) + |\nabla u|^2\diver X + \diver\big\{2X(u)\nabla u-|\nabla u|^2X\big\} \label{doga2} \\
\hphantom{2\langle\nabla_{\nabla u}X,\nabla u\rangle}{}= \bigl(|\nabla u|^2 - \varphi u^2\bigr)\diver X - X(\varphi)u^2 \notag\\
 \hphantom{2\langle\nabla_{\nabla u}X,\nabla u\rangle=}{}+ \diver\big\{2X(u)\nabla u+\bigl(\varphi u^2-|\nabla u|^2\bigr)X\big\}. \label{doga3}
\end{gather}
\end{Corollary}

\begin{proof}
If $\Delta u=\varphi u$, then
\begin{align*}
 2X(u)\Delta u &= 2\varphi X(u)u = \varphi X\bigl(u^2\bigr) = X\bigl(\varphi u^2\bigr) - X(\varphi)u^2 \\
 &= \diver\bigl(\varphi u^2X\bigr) - \varphi u^2\diver X - X(\varphi)u^2.
\end{align*}
Together with \eqref{rell}, this gives \eqref{doga2} and \eqref{doga3}.
\end{proof}

\setcounter{section}{3}
% Original source lines 492--585
\section{Absolutely continuous spectrum}%\label{secxavier}
%%%%%%%%%%%%%%%%%%%%%%%%%%%%%%%%%%%%%%%%%%%%%%

Let $S$ be a self-adjoint operator in a separable Hilbert space $H$,
and denote by $R(z)=(S-z)^{-1}$ the resolvent of $S$.
Say that a closed operator $T$ in $H$ is \emph{$S$-smooth} if, for each $x\in H$ and $\ve\ne0$,
$R(\lambda+{\rm i}\ve)x\in\dom T$
 for almost all $\lambda\in\R$ and
\begin{align*}%\label{Kato1}
 \sup_{|x|=1,\,\ve\ne0} \int_{-\infty}^\infty |TR(\lambda+{\rm i}\ve)x|^2 < \infty,
\end{align*}
see~\cite[p.~142]{ReedSimon78}.
In our context, the point is that then the image
\begin{align}\label{Kato3}
 \ran T^* \subseteq H_{\ac}(S),
\end{align}
see~\cite[Theorem XIII.23]{ReedSimon78}.
For a Borel subset $B\subseteq\R$,
say that a closed operator $T$ in $H$ is \emph{$S$-smooth on $B$} if $TE_B$ is $H$-smooth,
where $E_B$ denotes the spectral projection of $S$ associated to $B$.
Then
\begin{align*}%\label{Kato5}
 \ran E_BT^* \subseteq H_{\ac}(S),
\end{align*}
by \eqref{Kato3}.
By~\cite[Theorem XIII.30]{ReedSimon78}, $T$ is $S$-smooth on the closure $\bar B$ of $B\subseteq\R$ if
\begin{align}\label{Kato7}
 \dom T\supseteq\dom S
 \qquad\text{and}\qquad
 \sup_{\substack{|x|=1, \, \lambda\in B,\\ 0<|\ve|<1}} |\ve||TR(\lambda+{\rm i}\ve)x| < \infty.
\end{align}
Xavier~\cite[p.~581f.]{Xavier1988} shows that the latter criterion is satisfied
if $B$ is bounded and if there is a~constant $C$ such that
\begin{align}\label{Kato9}
 |Tx|^2 \le C|(S-\lambda)x|(|Sx|+|x|)
\end{align}
for all $x$ in a core of $S$ in $H$ and all $\lambda\in B$.
In fact, he shows that the term in \eqref{Kato7} is then bounded by $C(1+|\lambda|+|\ve|)$,
which explains why $B$ is assumed to be bounded; see~\cite[top of~p.~582]{Xavier1988}.

Guided by~\cite[Section 3]{Xavier1988}, we return to the geometric situation in Section~\ref{secrell}
and let $X$ be a~$C^1$ vector field and $u$ be a $C^2$ function on a Riemannian manifold $M$
such that $\supp X\cap\supp u$ is compact.
Then \eqref{rell} gives, for any smooth domain $D$ in $M$ with outer normal $\nu$ along $\partial D$,
\begin{gather*}
 \int_D\langle X,\nabla u\rangle\Delta u + \frac12\int_D|\nabla u|^2\diver X \\
 \qquad\quad{}= \int_D\langle\nabla_{\nabla u}X,\nabla u\rangle
 - \int_{\partial D}\biggl(\langle X,\nabla u\rangle\langle\nabla u,\nu\rangle
 - \frac12|\nabla u|^2\langle X,\nu\rangle\biggr).
\end{gather*}
As in~\cite[p.~583]{Xavier1988}, we use
\begin{align*}
 |\nabla u|^2 = u\Delta u - \frac12\Delta u^2
\end{align*}
and obtain
\begin{align*}
 \int_D \Delta u\biggl(\langle X,\nabla u\rangle + \frac{u}2\diver X\biggr)
 = {}&\int_D \biggl(\langle\nabla_{\nabla u}X,\nabla u\rangle + \frac14\diver X\Delta u^2\biggr) \\
 &{}- \int_{\partial D}\biggl(\langle X,\nabla u\rangle\langle\nabla u,\nu\rangle
 - \frac12|\nabla u|^2\langle X,\nu\rangle\biggr).
\end{align*}
Since
\begin{align*}
 P_Xu = \langle X,\nabla u\rangle + \frac{u}2\diver X
\end{align*}
satisfies
\begin{align*}
 u P_Xu = \diver\biggl(\frac12u^2X\biggr),
\end{align*}
we arrive at
\begin{align}
 \int_D (\Delta u-\lambda u) P_Xu
 = &\int_D \biggl(\langle\nabla_{\nabla u}X,\nabla u\rangle + \frac14\diver X\Delta u^2\biggr) \notag\\
 &- \int_{\partial D}\biggl(\langle X,\nabla u\rangle\langle\nabla u,\nu\rangle
 - \frac12\bigl(|\nabla u|^2 - \lambda u^2\bigr)\langle X,\nu\rangle\biggr),\label{x1}
\end{align}
which corresponds to~\cite[identity~(2)]{Xavier1988}, but with boundary integral included.
It will also be useful to have the latter formula with one of the substitutions
\begin{align}
 \int_D\diver X\Delta u^2
 &= \int_D\big\langle\nabla\diver X,\nabla u^2\big\rangle - \int_{\partial D}\diver X\nabla_\nu u^2 \label{x2a}\\
 &= \int_D(\Delta\diver X)u^2
 + \int_{\partial D} \bigl((\nabla_\nu\diver X)u^2 - \diver X\nabla_\nu u^2\bigr), \nonumber%\label{x2b}
\end{align}
which is Green's formula applied to the functions $\diver X$ and $u^2$.
Clearly,
\begin{gather}
 \int_D (\Delta u-\lambda u) P_Xu
	\le \|\Delta u-\lambda u\|_2\|P_Xu\|_2
	\le C\|\Delta u-\lambda u\|_2(\|\Delta u\|_2+\|u\|_2)\label{x3}
\end{gather}
if $X$ and $\diver X$ are bounded on $D$,
and $u$ satisfies a non-positive boundary condition.
This will be important in view of~\eqref{Kato9}.


% Original source lines 641--731
\begin{proof}[Proof of first part of Theorem~\ref{onlyacm}]
Let $b$ be a Busemann function associated to $\xi$ and ${X=\nabla b}$.
Then $|X|=1$ and $\diver X=h>0$.
Furthermore, $\nabla X$ is positive definite on the orthogonal complement of $X$,
that is, in the direction to the horospheres in $M$ with center $\xi$.

Let $Y$ be a smooth vector field on $M$ with compact support such that $Y\perp X$.
Then there are constants $C_1,C_2>0$ such that $\nabla X\ge C_1$ on the support of $Y$ and perpendicularly to $X$
and such that $|Y|^2\le C_2$.

Now $C^\infty_{\rm c}(M)$ is a core for $\Delta$.
Since $\diver X=h>0$ is constant, \eqref{x1} and \eqref{x2a} yield, for $u\in C^\infty_{\rm c}(M)$,
\begin{align*}
 \int_M (\Delta u-\lambda u) P_Xu
 &= \int_M \langle\nabla_{\nabla u}X,\nabla u\rangle \ge C_1 \int_M|\nabla^\bot u|^2 \\
 &\ge \frac{C_1}{C_2}\int_M|Yu|^2 = \frac{C_1}{C_2}\|T_Yu\|_2^2,
\end{align*}
where $\nabla^\bot u$ denotes the component of $\nabla u$ perpendicular to $X$.
Now \eqref{x3} applies and shows that $T_Y$ is $S$-smooth on any given bounded Borel subset $A\subseteq\R$.
Hence \eqref{Kato3} implies that $\ran E_AT_Y^*\subseteq H_{\ac}(\Delta)$,
for any bounded Borel subset $A\subseteq\R$.
Therefore, $\ran T_Y^*\subseteq H_{\ac}(\Delta)$.

Now suppose that $u\in L^2(D )$ is perpendicular to $H_{\ac}(\Delta)$.
Then $u$ is perpendicular to $\ran T_Y^*$, for all vector fields $Y$ as above.
But then $T_Yu=0$ in the sense of distributions, for all such vector fields $Y$.
This implies that $u$ is constant along horospheres, that is, levels of $u$ are unions of horospheres.
Now the preimage $b^{-1}(B)$ of any Borel subset $B\subseteq\R$ of positive measure has infinite measure.
Therefore, $u$ vanishes, by square-integrability.
\end{proof}


\begin{proof}[Proof of part of Theorem~\ref{onlyacn}]
By the definition of elementary groups of isometries of $M$, the gradient field $X$ of Busemann functions,
that is, the field of velocity vectors of unit speed geodesics emanating from $\xi$,
is invariant under $\Gamma$.
Notice that $X$ spans the normal space to the fibers of the Riemannian submersion $\pi$.
Now the arguments of the previous proof apply and show that $u\in L^2(N)$ is perpendicular to $H_{\ac}(N)$
if and only if $u$ is constant on the fibers of $\pi$.
There are now two cases:
The fibers of $\pi$ have infinite volume.
This is always the case when $\Gamma$ is of type \ref{ax}.
Then, as in the previous proof, the square-integrability of $u$ implies that $u=0$.

In the second case, $\Gamma$ is of type \ref{pa}.
Then the flow $(F_t)_{t\in\R}$ of $X$ induces diffeomorphisms between the fibers of $\pi=b$,
and the volume element of horospheres is multiplied by $\exp(ht)$ under $F_t$.
Hence, since the fibers of $b$ have finite volume, their volumes satisfy
\begin{align*}
	\exp(ht) \big|b^{-1}(s)\big| = \exp(hs) \big|b^{-1}(t)\big|.
\end{align*}
Furthermore, the space $L^2_b(N)$ of functions on $N$, which are constant fiberwise,
and its orthogonal complement $L^2_0(N)$ are invariant under $\Delta$.
That is, the spectrum of $\Delta$ is the combination of the spectra of $\Delta$ on $L^2_b(N)$
and of $\Delta$ on $L^2_0(N)$.
By what we said above, if $u\in L^2_0(N)$ is perpendicular to $H_{\ac}(N)$, then $u=0$.
Thus $L^2_0(N)\subseteq H_{\ac}(N)$, and we are left with discussing $\Delta$ on $L^2_b(N)$.

Note that $u\in L^2(N)$ is in $L^2_b(N)$ if and only if $u$ is a pull-back $b^*v$ for some function $v$ on $\R$.
We endow $\R$ with the measure $\mu=h_0\exp(ht){\rm d}t$,
where $h_0=\big|b^{-1}(0)\big|$ and get that
\begin{align*}%\label{unitary}
	b^* \colon \ L^2(\R,\mu) \to L^2_b(N)
\end{align*}
is a unitary transformation.
With respect to $b^*$, the Laplacian on $N$ corresponds to a diffusion operator on $\R$,
\begin{align*}%\label{diffusion}
	\Delta b^*v = b^*\bigl(- v'' - hv'\bigr).
\end{align*}
Let $\vf$ be the square root of $h_0\exp(ht)$.
Then
\begin{align*}%\label{diffusion2}
	\vf' = \frac{h}2\vf \qquad\text{and}\qquad \vf'' = \frac{h^2}4\vf.
\end{align*}
Therefore, under the unitary transformation
\begin{align*}
	L^2(\R,\mu) \to L^2(\R), \qquad v \mapsto \vf v,
\end{align*}
we obtain
\begin{align*}
	(\vf v)'' = \vf''v +2\vf'v' + \vf v'' = \vf\biggl(\frac{h^2}4 v + hv' + v''\biggr).
\end{align*}
We conclude that the Schr\"odinger operator $Sw=-w''+h^2w/4$ in $L^2(\R)$
corresponds to the above diffusion operator on $L^2(\R,\mu)$.
Since $h^2/4$ is a constant, the spectrum of the latter is absolutely continuous and equal to $\bigl[h^2/4,\infty\bigr)$.
By unitary equivalence, the same holds for $\Delta$ on~$L^2_b(N)$.

By Observation~\ref{chee}, $\sigma(N)\subseteq\bigl[h^2/4,\infty\bigr)$, hence the above shows equality,
$\sigma(N)=\bigl[h^2/4,\infty\bigr)$, in the case where the fibers of $\pi$ are of finite volume.
The case where the fibers are of infinite volume will be discussed in Section~\ref{secsig}.
\end{proof}

\setcounter{section}{5}
% Original source lines 920--1047
\section{Determination of the spectrum}\label{secsig}
%%%%%%%%%%%%%%%%%%%%%%%%%%%%%%%%%%%%%%%%%%%%%%

In this section we aim at the determination of the spectrum $\sigma(M)$.
We start with the case where $M$ is a Hadamard manifold with sectional curvature $K_M<0$
which is asymptotically harmonic about a point $\xi\in M_\infty$
with mean curvature of horospheres equal to $h>0$.
Recall from the first part of the proof of Theorem~\ref{onlyacn}
that the case left is $N=\Gamma\backslash M$,
where the fibers of the associated Riemannian submersion $\pi=b\colon N\to\R$ are of infinite volume,
where $b$ is the push-down to $N$ of some Busemann function associated to $\xi$.
We subsume the case $N=M$ here by including $\Gamma=\{\id\}$.
From Observation~\ref{chee}, we know that $\sigma(N)\subseteq\bigl[h^2/4,\infty\bigr)$.

For any smooth function $v$ on $\R$, we have
\begin{align*}
	\Delta b^*v = b^*\bigl(-v'' - hv'\bigr).
\end{align*}
Now $\vf=\vf(t)=\exp(ht/2)$ satisfies $\vf'=h\vf/2$ and $\vf''=h^2\vf/4$.
With $\mu=\vf^2{\rm d}t$, the unitary transformation
\begin{align*}
	\Phi \colon \ L^2(\R,\mu) \to L^2(\R), \qquad \Phi(v) = \vf v,
\end{align*}
yields
\begin{align*}
	(\Phi(v))'' = (\vf v)'' = \vf\biggl(\frac{h^2}4v + hv' + v''\biggr),
\end{align*}
and hence $\Phi$ intertwines the diffusion operator $Lv=-(v''+hv')$ in $L^2(\R,\mu)$
with the Schr\"odinger operator $Sw=-w''+h^2w/4$ in $L^2(\R)$.
Therefore, the spectra of $L$ and $S$ coincide,
where we recall that $H_{\ac}(S)=L^2(\R)$ with $\sigma(S)=\bigl[h^2/4,\infty\bigr)$.

We want to show that $\sigma(S)\subseteq\sigma(N)$.
In contrast to the case where the volume of the fibers of $b$ are of finite volume,
we cannot simply use pull backs $b^*v$ since we would lose square-integrability that way.
What we will do is to carefully cut off pull-backs.

\begin{Lemma}
Suppose that $-1\le K_M\le-a^2<0$ and that $\|\nabla R_M\|_\infty<\infty$.
Let $\chi_0\colon b^{-1}(0)\to\R$ be a non-zero $C^2$ function with compact support.
Extend $\chi_0$ to $N$ by letting $\chi=\chi_t$ be equal to~$F_{-t}^*\chi_0$ on $b^{-1}(t)$,
where $(F_t)_{t\in\R}$ denotes the flow of $X=\nabla b$.
Then $\chi$ is $C^2$ and $\nabla\chi$ and $\Delta\chi$ tend to zero uniformly as $t\to\infty$.
\end{Lemma}

\begin{proof}
Since $\chi$ is constant along the flow lines of $X$,
the gradient of $\chi$ is tangential to the fibers of $b^{-1}(t)$ of $b$.
By the upper curvature bound $-a^2$, unstable Jacobi field $J$ grow at least exponentially,
\begin{align*}
	|J(t)| \ge {\rm e}^{at} |J(0)| \qquad\text{for all} \quad t>0.
\end{align*}
Therefore,
\begin{align*}%\label{enab}
	|\nabla\chi_t| \le {\rm e}^{-at} |\nabla\chi_0| \qquad\text{for all}\quad t>0,
\end{align*}
which implies the first assertion.

For the estimate of the Laplacian, we refer to~\cite[Section 6]{BP21b}.
The situation there is that of a Hadamard manifold $X$ with pinched negative sectional curvature
and uniformly bounded covariant derivative of its curvature tensor,
a convex domain $C$ in $X$,
and a function obtained by extending a given function on $\partial C$
to $X\setminus C$ constantly along minimizing geodesics to $C$.
In our situation, $X$ corresponds to $M$, $C$ to the horoball $b^{-1}((-\infty,0])$,
and the given function to $\chi_0$.
A short look at the arguments in~\cite[Section 6]{BP21b} shows
that they also apply in the corresponding situation in $N$.
The estimate obtained in~\cite{BP21b} is that
\begin{align*}
	|\Delta\chi| \le c_0{\rm e}^{-at} \qquad\text{over} \quad b^{-1}(t)\quad \text{for all} \quad t>0,
\end{align*}
where $c_0$ is a constant, which depends on bounds for $K_M$, $\nabla R_M$, $\nabla\chi_0$, and $\nabla^2\chi_0$;
compare with~\cite[formula~(6.4) and end of Section~6]{BP21b}.
\end{proof}

\begin{proof}[End of proof of Theorem~\ref{onlyacn}]
Let $\lambda\in\sigma(S)$ and $\ve>0$.
Choose a non-vanishing $w\in C^\infty_{\rm c}(\R)$ such that $\|(S-\lambda)w\|_2\le\ve\|w\|_2$.
By shifting $w$ along $\R$ if necessary, we can assume that $\supp w\subseteq[t,\infty)$,
where $t>0$ is such that $|\Delta\chi| < \ve$ on $b^{-1}([t,+\infty))$.
Note that shifting $w$ is legitimate since $Sw$ is only shifted accordingly.
Then $v=w/\vf$ also has support in $[t,\infty)$ and satisfies $\|(L-\lambda)v\|_2\le\ve\|v\|_2$,
where now the $L^2$-norm is taken with respect to $\mu=\vf^2{\rm d}t$.
Then
\begin{align*}
	\Delta(\chi b^*v)
	= (\Delta\chi)b^*v - 2\langle\nabla\chi,\nabla b^*v\rangle + \chi\Delta b^*v
	= (\Delta\chi)b^*v + \chi b^*Lv
\end{align*}
since $\nabla\chi$ is perpendicular to $X$ and $\nabla b^*v$ is collinear with $X$
and since $b^*$ intertwines $\Delta$ with $L$.

Since $b$ is a Riemannian submersion and the flow maps $F_s$ of $X$ induces
a diffeomorphism of~$b^{-1}(0)$ with $b^{-1}(s)$,
for all $s\in\R$ and with respective Jacobian $\vf^2(s)={\rm e}^{hs}$.
Thus we identify~$x\in N$ with $(y,s)\in b^{-1}(0)\times\R$, where $b(x)=s$ and $F_s(y)=x$,
and get
\begin{align*}%\label{n1}
	\|\chi b^*v\|_2^2
	&= \int \chi^2(x)(b^*v(x))^2{\rm d}x \\
	&= \iint_{b^{-1}(0)\times\R}\chi_0^2(y)v^2(s)\vf^2(s){\rm d}y{\rm d}s \\
	&= \|\chi_0\|_2^2 \int_\R v^2(s)\vf^2(s){\rm d}s
	= \|\chi_0\|_2^2 \|v\|_2^2,
\end{align*}
where $\|\chi_0\|_2$ denotes the $L^2$-norm of $\chi_0$ in $L^2\bigl(b^{-1}(0)\bigr)$.
Likewise, by the choice of $t$,
\begin{align*}%\label{n2}
	\|(\Delta\chi)b^*v\|_2^2
	\le \ve^2|\supp\chi_0| \|v\|_2^2.
\end{align*}
and hence
\begin{align*}%\label{n3}
	\|(\Delta\chi)b^*v\|_2^2
	\le \ve^2\frac{|\supp\chi_0|}{ \|\chi_0\|_2^2}\|\chi b^*v\|_2^2
	= \ve^2c_1^2\|\chi b^*v\|_2^2,
\end{align*}
where $c_1$ depends on the choice of $\chi_0$.
Hence we obtain
\begin{align*}%\label{n4}
	\|(\Delta-\lambda)(\chi b^*v)\|_2^2
	&\le 2\|\chi b^*(L-\lambda)v\|_2^2 + 2\ve^2c_1^2\|\chi b^*v\|_2^2 \\
	&= 2\|\chi_0\|_2^2 \|(L-\lambda)v\|_2^2 + 2\ve^2c_1^2\|\chi b^*v\|_2^2 \\
	&= \ve^2c_2^2\|\chi b^*v\|_2^2,
\end{align*}
where $c_2$ depends on the choice of $\chi_0$.
This implies that $\lambda\in\sigma(N)$.
\end{proof}
\begin{thebibliography}{99}
\bibitem[3]{Ba95}
Ballmann W., Lectures on spaces of nonpositive curvature, \textit{DMV Seminar},
 Vol.~25, \href{https://doi.org/10.1007/978-3-0348-9240-7}{{Birkh\"auser}}, Basel, 1995.

\bibitem[4]{BMM17}
Ballmann W., Matthiesen H., Mondal S., Small eigenvalues of surfaces of finite
 type, \href{https://doi.org/10.1112/S0010437X17007291}{\textit{Compos. Math.}} \textbf{153} (2017), 1747--1768,
 \href{https://arxiv.org/abs/1506.06541}{arXiv:1506.06541}.

\bibitem[5]{BP21b}
Ballmann W., Polymerakis P., On the essential spectrum of differential
 operators over geometrically finite orbifolds, \textit{J.~Differential Geom.},
 {t}o appear, \href{https://arxiv.org/abs/2103.13704}{arXiv:2103.13704}.

\bibitem[6]{BLM10}
Bessa G.P., Jorge L.P., Montenegro J.F., The spectrum of the
 {M}artin--{M}orales--{N}adirashvili minimal surfaces is discrete,
 \href{https://doi.org/10.1007/s12220-009-9101-z}{\textit{J.~Geom. Anal.}} \textbf{20} (2010), 63--71, \href{https://arxiv.org/abs/0809.1173}{arXiv:0809.1173}.

\bibitem[7]{DoGa92}
Donnelly H., Garofalo N., Riemannian manifolds whose {L}aplacians have purely
 continuous spectrum, \href{https://doi.org/10.1007/BF01444709}{\textit{Math. Ann.}} \textbf{293} (1992), 143--161.

\bibitem[8]{DoGa97}
Donnelly H., Garofalo N., Schr\"odinger operators on manifolds, essential
 self-adjointness, and absence of eigenvalues, \href{https://doi.org/10.1007/BF02921722}{\textit{J.~Geom. Anal.}}
 \textbf{7} (1997), 241--257.

\bibitem[13]{KazdanWarner74c}
Kazdan J.L., Warner F.W., Curvature functions for compact {$2$}-manifolds,
 \href{https://doi.org/10.2307/1971012}{\textit{Ann. of Math.}} \textbf{99} (1974), 14--47.

\bibitem[16]{ReedSimon78}
Reed M., Simon B., Methods of modern mathematical physics.~{IV}. {A}nalysis of
 operators, Academic Press, New York, 1978.

\bibitem[17]{Rellich43}
Rellich F., \"{U}ber das asymptotische {V}erhalten der {L}\"osungen von
 {$\Delta u+\lambda u=0$} in unendlichen {G}ebieten, \textit{Jahresber. Dtsch.
 Math.-Ver.} \textbf{53} (1943), 57--65.

\bibitem[19]{Xavier1988}
Xavier F., Convexity and absolute continuity of the {L}aplace--{B}eltrami
 operator, \href{https://doi.org/10.1007/BF01462884}{\textit{Math. Ann.}} \textbf{282} (1988), 579--585.

\end{thebibliography}
\end{document}
